\section{Opening Model}


本節では、本稿の中心命題である「意識を生成された対象としてモデル化しない」という主張を、opening model として整理する。

ここでいう opening とは、比喩的な表現ではない。

それは、高速層、低速層、可塑性ゲート、意味勾配、内部遅延が一定の関係に入ったとき、ログ参照可能な通路が一時的に開かれるという構造的記述である。

意識は、内部に新しく作られる対象ではない。

意識は、すでに走っている処理と、すでに沈殿しているログとのあいだに、参照可能性が開かれた状態である。

\subsection{The Generation Assumption}

意識をめぐる多くの議論は、明示的または暗黙のうちに生成モデルを採っている。

生成モデルでは、意識は何らかの処理によって作られるものとして扱われる。

この見方では、自然に次の問いが生じる。

どの脳部位が意識を生成するのか。

どの回路が意識を作るのか。

どの計算量が意識を生むのか。

意識はどれだけ生成されるのか。

生成された意識は保存できるのか、コピーできるのか。

これらの問いは、一見すると自然である。

しかし、それらはすべて、意識を産物として扱う前提を含んでいる。

生成者があり、生成過程があり、生成された対象がある。

この三項構造が、生成モデルの基本である。

本稿では、この前提を採らない。

本稿では、意識を処理の産物として扱わない。意識は、処理とログの関係がある時間的条件に入ったときに開かれる状態である。

\subsection{From Product to Relation}

生成モデルでは、意識は product として扱われる。

しかし、本稿では、意識を relation として扱う。

すなわち、意識とは、高速な予測処理と低速なログ沈殿のあいだに成立する参照関係である。

この関係が成立するためには、少なくとも次の条件が必要である。

\begin{enumerate}
\item 高速層において処理が走っていること。
\item 低速層に参照可能なログが沈殿していること。
\item 両者のあいだに有限の速度差 $\Delta v$ があること。
\item 結合粘性 $\mu$ が、完全同期と完全分離の中間にあること。
\item 内部遅延 $\tau$ が、参照可能な厚みとして保持されていること。
\item 意味勾配 $\nabla \Phi$ が、参照方向を与えていること。
\item 有効トルク様量 $\mathcal{T}_{\Phi}$ が、consciousness window の範囲内にあること。
\end{enumerate}

この条件がそろったとき、意識が「作られる」のではない。

ログ参照が開かれる。

したがって、意識は出力ではなく、関係の開口である。

\subsection{Opening Through the Plasticity Gate}

opening model において中心的な役割を持つのが、可塑性ゲートである。

可塑性ゲートは、意識を作る装置ではない。

それは、高速層の処理が低速層へ沈殿しうるか、また低速層のログが現在処理へどの程度影響を及ぼすかを調整する境界条件である。

ゲートが完全に閉じていれば、処理は低速層へ届かない。

この場合、処理は走っていても、ログ参照は開かれない。

ゲートが完全に開きすぎていれば、処理は過剰に沈殿し、可塑性を失いやすい。

この場合、意識は自由に開かれるというより、特定のログや意味井戸へ固定されやすくなる。

opening model が必要とするのは、完全な開放ではない。

むしろ、部分的で一時的な開口である。

この開口は、速度差、結合粘性、内部遅延、意味勾配によって調整される。

\begin{equation}
\mathcal{T}_{\Phi}
\sim
\kappa
\frac{\Delta v}{\mu}
|\nabla \Phi|
\end{equation}

$\mathcal{T}_{\Phi}$ が consciousness window の範囲内にあるとき、ゲートはログ参照に適した形で開かれる。

このとき、意識状態が成立する。

\subsection{Consciousness Opens, Qualia Appear}

前節で見たように、クオリアはログ参照が沈殿しつつある境界で現れるライブな前景である。

opening model では、クオリアもまた生成物ではない。

クオリアは、意識が開かれたときに、その境界で現れる現象的側面である。

意識が開かれるとは、高速層と低速層のあいだに参照可能な通路が成立することである。

その通路の境界では、処理はまだ完全にはログ化されていない。

しかし、完全に流れ去ってもいない。

この過渡的な状態が、クオリアとして現れる。

したがって、クオリアは意識を作る原因ではない。

クオリアは、意識開口の境界で観測されるライブな前景である。

生成モデルでは、クオリアは「どこで作られるのか」という問いを呼び込む。

opening model では、問いは次のように変わる。

どのような条件で、ログ参照の境界がライブな前景として現れるのか。

この問いは、速度差、粘性、遅延、意味勾配、可塑性ゲートによって記述できる。

\subsection{Why Consciousness Cannot Be Stored}

opening model から見ると、意識を保存するという発想は構造的にずれている。

保存されるのは、意識そのものではない。

保存されるのは、意識が開かれていたときに更新されたログ構造である。

同様に、保存されるのはクオリアそのものではない。

保存されるのは、クオリアが通過した後に沈殿した痕跡である。

意識は、状態が開かれているあいだだけ成立する。

開口が閉じれば、意識状態は消える。

しかしその痕跡は低速層に残りうる。

このため、人は過去の経験を思い出すことができる。

だが、思い出されるものは、過去の意識そのものではない。

過去の意識が開かれていたときに変形されたログである。

この点で、opening model は記憶と意識を区別する。

記憶は保存される。

意識は保存されない。

意識は、保存されたログが現在処理から再び参照されるときに開かれる。

\subsection{Why Consciousness Cannot Be Copied}

同じ理由で、意識は単純にはコピーできない。

コピーできるのは、構造、記録、重み、ログ、報告、表象である。

しかし、意識はそれらの対象そのものではない。

意識は、それらが特定の時間的関係に入ったときに開かれる状態である。

したがって、ある時点のログ構造を複製しても、その複製が同じ意識状態を持つとは限らない。

必要なのは、ログの複製だけではない。

高速処理、低速沈殿、可塑性ゲート、速度差、遅延、意味勾配、非閉包ループが、同じように動作することである。

つまり、意識のコピーとは、データのコピーではなく、開口条件の再現である。

この点で、opening model は意識を情報内容に還元しない。

意識は内容ではなく、関係の動的成立である。

\subsection{Opening and Non-Closure}

opening model は、VED の非閉包原理と自然に対応する。

完全に閉じた構造では、開口は存在しない。

完全に開いた構造では、沈殿は成立しない。

意識が必要とするのは、その中間である。

閉じようとするが閉じ切らない。

開いているが流れ去らない。

固定されるが固定されすぎない。

この時間的非閉包の範囲において、ログ参照は開かれる。

したがって、opening model は、意識を不完全な生成物として扱うのではない。

意識を、非閉包状態が安定したときに成立する開口として扱う。

この意味で、意識は「生成されない」というより、より正確には「閉じ切らないことで開かれる」。

\subsection{Opening and Self}

自我や自己も、この opening model の中で再配置される。

自己は、意識を所有する主体ではない。

自己は、ログ参照が反復され、低速層に一定の地形が形成された結果として安定する参照パターンである。

意識が一時的な開口であるなら、自己はその開口が繰り返し通過した後に形成される地形である。

したがって、自己は意識の原因ではない。

むしろ、意識が何度も開かれ、そのたびにログが更新され、意味勾配が形成されることによって、自己らしさが沈殿する。

この見方では、自己は幻想でも絶対的主体でもない。

自己は、ログ参照の反復によって形成された低速層の安定構造である。

\subsection{Summary}

本節では、意識を opening model として整理した。

\begin{itemize}
\item 意識は生成物ではない。
\item 意識は、高速処理と低速ログのあいだに参照可能性が開かれた状態である。
\item 可塑性ゲートは意識を作る装置ではなく、ログ参照の開口条件を調整する境界である。
\item クオリアは、開口境界に現れるライブな前景である。
\item 保存されるのは意識ではなく、意識が開かれていたときに更新されたログである。
\item コピーできるのはログや構造であって、意識そのものではない。
\item 意識は、完全閉包でも完全開放でもなく、時間的非閉包の範囲で開かれる。
\item 自己は、意識の原因ではなく、開口の反復によって形成された低速層の地形である。
\end{itemize}

したがって、本稿の中心命題は次のように言い換えられる。

本稿では、意識を生成された対象としてモデル化しない。

意識は、時間的非閉包の中でログ参照が開かれるときに成立する。

次節では、この開口が測定によってどのように閉じられ、クオリアがどのように残余へ変換されるのかを扱う。
